\section{CW-Complejos y tipo de homotopía}
%tipo de homotopia-----------------------
\begin{definition}
Sea $X$ un espacio topológico y $A \subset X$ un subespacio. Decimos que $A$ es un
\textit{retracto de deformación fuerte} de $X$ si existe una homotopía relativa a $A$ de la función identidad de $X$, $id_X$, a una retracción $r: X \to A$.
\end{definition}

\bigskip

\begin{definition}
Sean $X$ e $Y$ espacios topológicos. Decimos que $X$ e $Y$ son
\textit{homotópicamente equivalentes} si existen aplicaciones continuas
\begin{equation*}
    f: X \to Y, \qquad g: Y \to X
\end{equation*}
tales que
\begin{equation*}
    g \circ f \simeq \mathrm{id}_X \qquad \text{y} \qquad f \circ g \simeq \mathrm{id}_Y.
\end{equation*}
En este caso, decimos que $X$ e $Y$ tienen el mismo \textit{tipo de homotopía} y llamamos a $f$ y $g$ \textit{equivalencias homotópicas}.
\end{definition}

\medskip

\begin{nota}
    En lo que sigue nos referiremos a los retractos de deformación fuerte simplemente como \textit{retractos de deformación}.
\end{nota}

\begin{remark}
    Nótese que todo retracto de deformación de un espacio tiene el mismo tipo de homotopía que este. Dados un espacio $X$ y un retracto de deformación $A\subseteq X$, la función inclusión, $\iota: A\hookrightarrow X$, y la retración $r: X \to A$ son equivalencias homotópicas.
\end{remark}
%--------------------------------------------------------
\bigskip
%CW-complejos-------------------------------
\begin{definition}
    Una \textit{célula} \textit{e} es un espacio topológico homeomorfo al disco unidad de n-dimensional.Diremos que la dimensión de la célula es \textit{n} y diremos que \textit{e} es una \textit{n-célula}.
\end{definition}

\bigskip

\begin{definition}
    Llamamos \textit{CW-complejo} a un espacio topológico construido siguiendo los siguientes pasos:

        \begin{enumerate}[label=(\roman*)]
            \item Se comienza con un conjunto discreto $ X_0 $ (llamado también \textit{0-esqueleto}), cuyos elementos llamaremos \textit{0-células}.
            
            \item Se construye, de manera inductiva, el llamado $ n $-esqueleto $ X_n $ (donde $n$ denota la dimensión de las células que lo forman) a partir del $(n-1)$-esqueleto. 
            
            Tomando $ X_{n-1} $, se le pega una colección $ \Lambda $ de $ n $-células mediante aplicaciones $ \varphi_\lambda $ del $ (n-1) $-disco unidad a $ X_{n-1} $. 
            
            El espacio $ X_n $ es el espacio cociente de la unión disjunta $X_{n-1}\sqcup_{\lambda} e_\lambda$ por las identificaciones naturales asociadas a las aplicaciones $ \varphi_\lambda $ anteriores.
            
            \item Puede pararse tras un número finito de pasos (aunque no tiene por qué hacerse). Llamaremos \textit{CW-complejo} al espacio 
            \begin{equation*}
                X = \bigcup_{n} X_n.    
            \end{equation*}
            Si el número de pasos es infinito, los conjuntos abiertos de la topología serán aquellos que sean abiertos para cada $ X_n $.
        \end{enumerate} 
\end{definition}

{\color{blue}¿Pongo ejemplos de CW-complejos?}

%---------------------------------------------------------------